\chapter{CNA and XNA 4.0}
\label{ch:cna-xna}
\label{ch:what-is-cna}
\index{XNA}\index{API compatibility}\index{differential testing}

\section{The one-sentence version}

CNA is a native C\texttt{++} reimplementation of the \emph{XNA 4.0} programming model, compiled
with a C\texttt{++}23 toolchain floor but written in a disciplined C\texttt{++}17/20 idiom. It
reimplements the game framework Microsoft shipped for the Xbox 360 and Windows from roughly 2006
to 2013. CNA's default platform and audio implementations use SDL3, but those are independent
configuration axes. Rendering is delegated to one of \RendererIdentityCount{} public identities
mapped to \RendererFamilyCount{} implementation families; an opt-in build may compile several
families and select one at runtime.

That description contains three separate engineering problems which organize the rest of the
book:

CNA is a \emph{translation target}, not a loader or binary compatibility layer.  It cannot run
an existing C\# assembly.  Porting translates the source to C\texttt{++}; even a property access
becomes a \texttt{getXProperty()} or \texttt{setXProperty(value)} call.  Names and observable
contracts can remain XNA-shaped while the source language and ownership model cannot.

\begin{enumerate}
  \item \textbf{Reproduce an API.} XNA's public surface --- \cnans{Microsoft.Xna.Framework}
        and its sub-namespaces --- has to exist in C\texttt{++}, with the same class names,
        method signatures, namespaces, and observable behavior, so that code written against
        real XNA (or against FNA, the open-source XNA reimplementation in C\# that this
        project treats as its primary behavioral reference) has somewhere to go.
  \item \textbf{Reproduce a runtime.} XNA is a managed-code API. It assumes things a bare
        C\texttt{++} program does not have for free: garbage-collected \texttt{System.Object},
        events and delegates, exceptions with a .NET-shaped hierarchy, \texttt{TimeSpan},
        collections. CNA does not reimplement the CLR --- instead it depends on a sibling
        project, \texttt{sharp-runtime}, which provides a deliberately partial C\texttt{++}
        subset of the .NET base class library, just deep enough to make the XNA API's
        vocabulary make sense in C\texttt{++}.
  \item \textbf{Reproduce a GPU.} XNA only had to run against Direct3D 9 (Xbox 360's
        own GPU API is Direct3D-shaped too). CNA has to make the same \cnaclass{GraphicsDevice}
        /\cnaclass{SpriteBatch} / \cnaclass{Effect} surface available while pixels are produced
        by SDL's 2D renderer, OpenGL, Vulkan, BGFX, WebGPU,
        Direct3D 9/11/12 under Wine, an HTML \texttt{<canvas>} under Emscripten, or a
        from-scratch DirectDraw reimplementation. Identical API presence does not imply identical
        behavior or evidence; Part~\ref{part:renderers} records those distinctions.
\end{enumerate}

\begin{sourcenote}
\textbf{From the project's own README (2026-07-11 snapshot):} ``CNA is a C\texttt{++}
reimplementation of the XNA 4.0 programming model, built on SDL3 and a pluggable graphics
renderer layer. It is a framework/runtime and abstraction layer---not a game---designed to
preserve XNA-style APIs (\texttt{Microsoft::Xna::Framework}) while using modern C\texttt{++}
internals.''
\end{sourcenote}

\section{Why XNA, and why now}

XNA Game Studio was Microsoft's managed-code (C\#) game development framework, built around
a small, coherent \cnaclass{Game} / \cnaclass{GraphicsDevice} / \cnaclass{SpriteBatch} / \cnaclass{Effect}
vocabulary that a whole generation of independent and student developers learned game
programming on. Microsoft discontinued active XNA development years ago; the community
response was \textbf{FNA}, a faithful open-source reimplementation of the XNA 4.0 API in
C\#. The pinned local reference selects its SDL3 platform implementation by default and retains SDL2 as
an explicit compatibility choice through \texttt{FNA\_PLATFORM\_BACKEND=SDL2}; describing FNA
simply as an SDL2 runtime is therefore historical shorthand, not its current default.

CNA asks a different question than FNA does. FNA asks: \emph{how do we keep running C\#
XNA code forever, on modern platforms, without Microsoft?} CNA asks: \emph{what if the code
itself did not have to be C\#?} A native C\texttt{++} implementation of the same API gives a
team that likes the XNA/MonoGame mental model --- a thin, explicit, immediate-mode-ish 2D/3D
API without an editor or a scene graph imposed on you --- a path that does not require a
managed runtime or GC pauses, and that a build system and toolchain of your choosing can
own end to end. CNA's own README states this directly: it is meant to ``provide a native
C\texttt{++} path for teams that like the XNA/MonoGame model but need non-managed
runtime/toolchain control.''

This is also, not incidentally, why CNA is useful as a target for \emph{porting} existing
XNA/FNA games rather than only for writing new ones from scratch --- a use case this book
returns to in this book's migration and case-study chapters, using the real ports of
\emph{Speedy Blupi} and \emph{Planet Blupi} (via the \texttt{free-direct} sibling project)
as worked examples.

``Why now'' is not merely an argument from intent. At the pin, CNA has fifteen physical
framework modules, an optional native C API source module (whose alpha.1 final target is
compile-blocked), \RendererIdentityCount{} public renderer identities over
\RendererFamilyCount{} implementation families, an XNB
reader stack, its own CNJ format, a shared runtime/offline glTF importer, and implemented Media
types backed by file parsing, FFmpeg and SDL3\_mixer.  Those facts establish breadth; they do not
establish uniform parity.  Later chapters therefore report each route's behavior, host and
oracle separately instead of reviving a fragile percentage of ``types present.''

\section{Design goals, in the project's own words}

CNA's own README states six goals directly, worth quoting in full since every one of them
recurs as a concrete design decision somewhere later in this book, not just as an aspiration:
``recreate the XNA developer experience in native C\texttt{++}''; ``provide a native
C\texttt{++} path for teams that like the XNA/MonoGame model but need non-managed
runtime/toolchain control''; ``mirror core XNA namespaces and API patterns while implementing
them incrementally''; ``decouple gameplay-facing API from renderer implementation
details''; ``enable one high-level API surface across different rendering technologies''; and
``keep SDL/OpenGL/Vulkan-level concerns behind framework abstractions.'' The middle four of
these six are really one idea stated four ways --- the renderer-abstraction goal --- and
Chapter~\ref{ch:backend-architecture} is where that idea becomes a concrete
\cnaclass{IGraphicsRenderer} contract rather than a design principle; the first two are the
``why native C\texttt{++}, why not just use FNA'' argument already made above.

\section{Verification is a vector, not a leaderboard}

There is no honest one-number answer to ``how tested is CNA?''  A source tree can contain a test
definition without registering it in a selected build; CTest can register a command that the host
cannot execute; a renderer can return success without consuming effect state; and a pixel can be
correct on one translated runtime without proving another host.

The book therefore records a claim as a vector: public identity and implementation family,
configuration, execution host, driver or translation layer, oracle, and unavailable branches.
Examples deliberately occupy different tiers:

\begin{itemize}
  \item the \RendererIdentityCount{}-identity / \RendererFamilyCount{}-family registry is strong structural evidence, not rendering proof;
  \item the Direct3D~9 corpus compares \XnaOracleSceneCount{} scenes with real XNA output, but
        runs through Wine and
        DXVK rather than native period hardware;
  \item HTML\_DOM has HTTP-served Chromium execution and compositor samples, while a successful
        build of another Web identity proves less;
  \item recorded Vulkan and BGFX pass fractions belong to their named campaigns and configured
        populations, not to a universal current suite total.
\end{itemize}

Appendix~\ref{ch:feature-matrix} gives the bounded renderer matrix.  Chapters~\ref{ch:counting-discipline}
and~\ref{ch:oracles} define the counting and oracle vocabulary used throughout the edition.

\section{What CNA is not}

It is worth being precise about the negative space, because CNA's own documentation is:

\begin{itemize}
  \item \textbf{Not a game.} There is no bundled playable demo executable shipped as the
        project's flagship artifact --- the README states this is a deliberate priority of
        framework/runtime development over a shipping game demo. What ships instead is a
        test suite (thousands of GoogleTest cases across the unit and GPU pixel-test suites)
        and a set of minimal verification executables like \texttt{hello-triangle-sdl}.
  \item \textbf{Not an XNA content builder.} CNA has a substantial runtime XNB reader, but no
        MGCB-style import/process/write toolchain and no automatic reader bootstrap. A game must
        register the built-in readers before loading XNB; CNJ, loose files, and glTF are separate
        routes. Chapters~\ref{ch:xnb-container}--\ref{ch:cnj-format} define those boundaries.
  \item \textbf{Not universally compiled-shader compatible.} CNA has distinct answers to
        XNA effects: renderer-owned stock effects; compiled XNA/FNA D3D9 Effect Framework bytecode
        accepted by the public \cnaclass{Effect} constructor on qualified renderers; and
        renderer-specific custom source or program forms through \cnaclass{ShaderEffect}.
        \texttt{.fx} or HLSL source, DXBC and MGFX are not interchangeable with that compiled
        Effect Framework format. No one route works identically on all \RendererFamilyCount{} families;
        Chapter~\ref{ch:shader-fx-gap}
        separates those contracts.
  \item \textbf{Not binary-compatible with real Xbox Live.} The \cnaclass{GamerServices}
        namespace (Chapter~\ref{ch:gamerservices}) reimplements the public API shape with
        local/synthetic semantics, matching how FNA itself already handles this namespace ---
        there is no real Xbox Live behind it, by design.
  \item \textbf{Not a drop-in Microsoft DirectX runtime.} The \texttt{free-direct} sibling
        implements a game-driven subset of DirectDraw, DirectSound, and DirectPlay over SDL3 and
        free-api. Its default DirectPlay transport is in-process; ENet networking is opt-in and
        speaks to the same reimplementation, not Microsoft's historical service. Chapter~\ref{ch:freedirect-freeapi}
        gives the per-method status census.
\end{itemize}

\section{Auditability as a framework property}

The single most distinctive thing about how this project is run --- and the reason this book
can make specific, checkable claims instead of vague ones --- is its verification
methodology. Three mechanisms recur throughout the codebase and throughout this book:

\begin{description}
  \item[Differential testing against a real reference implementation.] \texttt{tools/fna-reference/}
        runs a genuine, live \texttt{FNA.dll} and compares CNA's behavior against it
        call-for-call, rather than against a written specification of what FNA is supposed to
        do. Disputed behavior that differential testing cannot settle is checked against real
        XNA 4.0 running in a Windows~7 virtual machine.
  \item[A compile-time purity check.] A dedicated build define, \texttt{CNA\_STRICT\_XNA\_API},
        turns the otherwise-inert \texttt{CNAEXT} marker macro into a real
        \texttt{[[deprecated]]} attribute, checked under \texttt{-Werror=deprecated-%
        declarations} by two small standalone CMake targets built specifically for this
        purpose. Because every CNA extension beyond the real XNA 4.0 API must be wrapped in the
        \texttt{CNAEXT} marker macro (Chapter~\ref{ch:design-philosophy} explains the
        convention, and the self-verifying pair of targets that checks it, in full), this gives
        the project a mechanical --- not just documented --- guarantee that the public
        \cnaclass{Microsoft::Xna::Framework} surface has not silently grown project-specific
        extensions.
  \item[Pixel-level oracle corpora, not just ``it compiles.''] The \texttt{D3D9} renderer in
        particular ships an oracle corpus that diffs CNA's own render output against
        \emph{the real XNA 4.0 runtime's own render} of the same scene.  The checked-in corpus
        contains \XnaOracleSceneCount{} declarative scenes and the same number of matching
        reference PNGs. Its registered D3D9 gate runs the full corpus through the absolute
        per-channel comparator at tolerance zero; that
        wiring is an enforceable contract, while a pass count still requires an identified run.
        Chapter~\ref{ch:direct3d-backends} covers what that verification pipeline (MinGW-w64
        cross-compilation, Wine, DXVK, a real GPU) actually looks like.
\end{description}

\begin{sourcenote}
\textbf{Why ``differential testing against FNA'' is not by itself enough --- a real finding.}
CNA's own \texttt{CHECKLIST.md} names FNA as authoritative for \emph{behavior} but the original
Microsoft XNA 4.0 reference assemblies as authoritative for \emph{API surface}, and explains why
with a genuine near-miss: FNA is itself a reimplementation, and its own \texttt{src/Media/Song.cs}
quietly omits three public XNA 4.0 members --- \texttt{Song.Album}, \texttt{Song.Artist}, and
\texttt{Song.Genre} --- along with \texttt{Song.ToString()}. Because CNA's \cnans{Media} namespace
(Chapter~\ref{ch:media-system}) was itself audited against FNA across eight consecutive
adversarial review passes, every one of those eight passes came back clean --- CNA had faithfully
reproduced FNA's own incompleteness, not merely failed to find a bug that was there to find. The
gap was only caught on a ninth pass, once the audit method itself changed to a member-level diff
against the real Microsoft reference XML rather than another comparison against FNA. The fix
(tracked as \texttt{MEDIA-174} through \texttt{MEDIA-180} in \texttt{plan\_media.md}) added
non-owning \cnaclass{Album}\texttt{*} / \cnaclass{Artist}\texttt{*} / \cnaclass{Genre}\texttt{*}
back-pointers to \cnaclass{Song},
populated by \cnaclass{MediaLibrary} and \texttt{nullptr} for a song outside any library ---
mirroring an ownership pattern the surrounding \cnans{Media} types already used. The methodological
lesson the project drew from this, stated directly in its own checklist, is the one this chapter
opened with: ``FNA doesn't have it either'' is not a justification, because that is exactly the
reasoning that let this specific gap survive eight reviews.
\end{sourcenote}

This matters for how you should read the rest of this book. Where a claim in the chapters
ahead is stated as a fact about CNA's behavior, it is because the project's own test suite,
documentation, or source comments state it as a verified fact, generally with a specific
mechanism (a CTest binary, a doc file, a task number) behind it --- and where the project's
own documentation flags something as a dated snapshot, an estimate, or a known open gap, this
book preserves that qualification rather than erasing it for a cleaner sentence.

\section{The first tagged release and the continuing phase workflow}

CNA developed continuously from 2025-02-22 without conventional releases while numbered phases
and tasks carried engineering state. On 2026-08-20 the annotated
\CnaEditionTag{} tag established the first Semantic Versioning release. The product version is
assembled from CMake's numeric \texttt{0.1.0} project version and the separate
\texttt{alpha.1} prerelease field. CMake prints the complete version during configuration,
\texttt{cmake/Version.cmake} generates \path{CNA/Version.hpp}, Doxygen uses the same product
version, and the changelog and release procedure name the tag spelling with its leading
\texttt{v}.

The leading zero and alpha qualifier matter. The release policy permits compatibility-breaking
changes before 1.0 when the minor or prerelease version advances; alpha.1 should not be read as a
stable 1.0-quality API promise. The native C ABI has its own experimental version, 0.7.0, described
in Appendix~\ref{ch:native-c-api}; product and ABI versions answer different questions.

SemVer did not replace the project's strictly incrementing \textbf{Phase} counter. Phases in
\texttt{CHECKLIST.md} and the per-subsystem \texttt{plan\_*.md} files still name bounded work,
split it into numbered tasks, and retain closure evidence. They are an engineering workflow;
tags identify publishable source states. Chapter~\ref{ch:project-practice} covers how the two now
coexist.

This phase-based, task-tracked structure supplies valuable historical evidence, but a closing
count remains a snapshot of one configuration.  This edition preserves such a number only with
its campaign and verdict channel; current source definitions, CTest registrations, executed
cases, and oracle scenes are never silently combined into one total.

\section{CNA's software license}

CNA is licensed under the Microsoft Public License (Ms-PL), the same permissive-but-not-MIT
license FNA itself uses --- a deliberate choice, not a coincidence: portions of CNA are
directly derived from or based on FNA's own source (most visibly the D3D9 renderer's vendored
Stock Effects HLSL sources, discussed in Part~\ref{part:renderers}), and matching FNA's license keeps that
derivation legally uncomplicated. Ms-PL permits commercial use, modification, and
redistribution, but --- unlike MIT or the 2-clause BSD license --- includes an explicit patent
grant limited to contributions and a reciprocal ``if you sue over patents, your license ends''
clause; it is not a copyleft license, so a game built against CNA is not obligated to publish its
own source.

\section{Reading map for the rest of this book}

Chapter~\ref{ch:ecosystem-map} draws the full dependency graph between CNA and its sibling
repositories. Chapter~\ref{ch:design-philosophy} is the ``house style'' chapter --- the
concrete rules (API mirroring, namespace discipline, the \texttt{CNAEXT} tag, the
\texttt{getXProperty} / \texttt{setXProperty} convention, the porting checklist) that every
other chapter's code examples silently assume. Part~\ref{part:architecture-framework} gets you from a clean checkout to a
running triangle and then through the core framework types. Parts~\ref{part:graphics-machine}
and~\ref{part:renderers} are the
graphics machine, first as an API and then as \RendererFamilyCount{} implementation factories serving \RendererIdentityCount{} public
renderer identities.
